\documentclass{article}
\usepackage{neurips_2024}
\usepackage{amsmath, amssymb}
\usepackage{graphicx}
\usepackage{hyperref}
\usepackage{booktabs}

\title{AAES-OS v1.0: A Constitutional Runtime Architecture for Governed Intelligence}

\author{
  Jon Halstead \\
  \texttt{zeronull1983@aaes-os.org}
  \And
  Dar-z Morris \\
  \texttt{darz@aaes-os.org}
}

\begin{document}

\maketitle

\begin{abstract}
We introduce AAES-OS v1.0, a constitutional runtime architecture for governed AI systems. AAES-OS unifies three mathematical layers---Wave Math (micro-level judgment dynamics), Continuity Failure Theory (macro-level consequence propagation), and Reconstruction Sufficiency (meta-level reconstructability)---into a deterministic reference implementation for enforcing constitutional constraints on model-mediated behavior. The system includes CAS 1.0, a formal object model; CRK-1, a constitutional runtime kernel; CTS, a conformance test suite; CDP-1, a reproducible continuity benchmark; and CEP, an experimental platform enabling independent replication. Version 1.0 separates architectural claims, implementation results, and research hypotheses so evidence can be evaluated without overclaiming. Architectural claims describe properties guaranteed by the specification, implementation results describe measured outcomes demonstrated by the AAES-OS v1.0 reference implementation, and research hypotheses identify broader claims that require validation across independent implementations or production deployments.
\end{abstract}

\section{Introduction}
Modern AI systems exhibit drift, nondeterminism, and opaque internal reasoning. These properties undermine safety, reproducibility, and scientific evaluation. Existing governance approaches rely on post-hoc analysis or heuristic constraints that cannot guarantee constitutional compliance.

AAES-OS addresses this gap by introducing a constitutional architecture: a system where governance is enforced at runtime through deterministic rules, invariant validation, and content-addressed receipts. The paper's primary contribution is architectural: it defines a constitutional operating model, a reference implementation, and an evidence model that others can evaluate, reproduce, and extend. It does not claim to solve AI governance universally.

Our contributions are:
\begin{itemize}
    \item Normative contribution: constitutional specification, governance model, and architectural contracts.
    \item Engineering contribution: CRK-1 runtime, ledger, receipts, cockpit, CTS, CDP-1, CEP, and tooling.
    \item Research contribution: hypotheses about continuity, reconstruction, and semantic boundaries evaluated through evidence and replication rather than assumed universal.
\end{itemize}

\subsection{Claim Taxonomy}
AAES-OS v1.0 uses an explicit evidence taxonomy. Architectural claims are properties guaranteed by the specification, such as layer separation, provenance model, constitutional contracts, and governance model. Implementation results are measured outcomes demonstrated by the AAES-OS v1.0 reference implementation, such as test results, replay determinism, freeze mechanics, drift measurements, and ledger integrity. Research hypotheses are broader claims about governed intelligence that require validation across multiple implementations or production deployments.

\section{Related Work}
We review prior work in constitutional AI, formal verification, deterministic runtimes, reproducibility in ML, and governance architectures. AAES-OS differs by providing a complete, executable, constitutional runtime rather than a policy layer or training-time constraint.

\section{Mathematical Foundations}
\subsection{Wave Math}
Wave Math defines micro-level judgment dynamics and local decision stability.

\subsection{Continuity Failure Theory}
Continuity Failure Theory models macro-level consequence propagation and drift under perturbation.

\subsection{Reconstruction Sufficiency}
Reconstruction Sufficiency ensures that all system states are reconstructable from receipts and evidence.

\section{Constitutional Architecture}
\subsection{Invariants}
The architecture defines a frozen set of invariants (K-$\infty$, K0--K15, K$\Omega$) that govern all transitions.

\subsection{Object Model}
CAS 1.0 defines objects including Identity, Decision, Outcome, Evidence, Interpretation, Receipt, DriftObservation, and KernelChallenge.

\subsection{Allowed Transitions}
All transitions are deterministic and constitutionally validated.

\subsection{Deterministic Behavior Rules}
No hidden state, no nondeterministic branches, no side-effects outside receipts.

\section{CRK-1 Runtime}
CRK-1 is the first deterministic constitutional runtime. It includes:
\begin{itemize}
    \item Constitutional boot
    \item Proof gate
    \item Capability gate
    \item Receipt generation
    \item Drift detection
    \item Deterministic execution engine
\end{itemize}

\section{CAS 1.0}
CAS 1.0 provides:
\begin{itemize}
    \item A formal object model
    \item A reference implementation
    \item A conformance test suite (CTS)
\end{itemize}

\section{CDP-1 Benchmark}
CDP-1 measures continuity drift under controlled perturbations. It includes a dataset, metrics, thresholds, reproduction scripts, and continuity graphs.

\section{CEP Platform}
CEP enables experiment execution, logging, deterministic replay, trace capture, and reproducibility.

\section{Evaluation Frame}
We frame evaluation around three systems questions:
\begin{enumerate}
    \item Can the architecture be implemented? This is demonstrated through the reference implementation, deterministic replay, freeze mechanics, conformance tests, and release gates.
    \item Can the architecture be independently verified? This is demonstrated through provenance, content-addressed receipts, COR/CAR/CAV-style registries, deterministic replay, and the conformance ecosystem.
    \item Can the architecture evolve without losing constitutional integrity? This is demonstrated through amendment protocols, governance receipts, canonical freeze, and founder-independent stewardship.
\end{enumerate}

This framing separates architectural feasibility, implementation evidence, and long-term governance claims. It also prevents measured reference-implementation results from being mistaken for universal guarantees about all future intelligent systems.

\section{Independent Replication}
AAES-OS v1.0 includes a complete replication package enabling external teams to run CAS, CRK-1, CTS, and CDP-1, publish results, and challenge conclusions.

Replayability means the reference implementation reproduces the same result from the same recorded evidence. Independent verifiability means a separate implementation reaches the same result using only the specification and recorded artifacts. AAES-OS treats these as distinct evidentiary thresholds.

\section{Discussion}
The primary contribution of AAES-OS v1.0 is architectural. The paper presents a constitutional runtime model with inspectable governance, reproducible evidence, and founder-independent stewardship. It does not claim that constitutional runtime governance alone solves model alignment, semantic truth, or all production governance problems. Instead, it shows that governance can be made observable, replayable, and auditable at the runtime boundary.

Specifications define intended behavior. Implementations demonstrate one realization. Evidence establishes what has been achieved. Replication determines confidence.

\section{Conclusion}
AAES-OS v1.0 presents a constitutional runtime architecture for governed AI systems. Its contribution is not a priority claim about being first, nor a universal claim about solving AI governance. It is a coherent constitutional model whose architectural properties are specified, whose implementation results are measured, and whose broader research hypotheses remain open to independent validation. Its scientific posture is evidence-first: guarantees are specified, implementations are tested, empirical claims are backed by artifacts, and confidence grows through independent replication.

\bibliographystyle{plain}
\bibliography{references}

\end{document}
